\section{Actual coefficient presentations supply formal centre powers}

Write the original coordinates as $(u,y,z)$, with $u$ tangent to a smooth
centre, $y$ normal to it inside the presentation carrier, and $z$ normal to
that carrier. This chapter spells out the coefficient-ideal step used in
the proof of Bierstone--Milman Theorem 9.4. We use literal power-series
coefficients, and do not identify an unnormalized derivative with a
coefficient: any factorial normalization must be included in the actual
coefficient data. No Hilbert--Samuel equality or transformed presentation
is concluded here.

\begin{lemma}[Whole formal centre powers from actual coefficient series]
\label{hir:actual-formal-centre-power-relative-coefficients}
\uses{hir:actual-formal-coordinate-power-coefficient-kernel}
\lean{Hironaka.HironakaActualCoefficientPresentationCentrePower.actualJoinedExponent_split,Hironaka.HironakaActualCoefficientPresentationCentrePower.actualJoinedExponent_centre_weight,Hironaka.HironakaActualCoefficientPresentationCentrePower.mem_actualPresentationCentre_power_iff}
\leanok
Let $F\in\mathbb C[[w,z]]$, where $w\in\mathbb C^d$, $z\in\mathbb C^r$,
and let $A\subset\{1,\ldots,d\}$. Write $F=\sum_\gamma F_\gamma(w)z^\gamma$
using its literal coefficients. Let $P_A$ be the actual ideal of the
selected $w$ variables and let $P_{A,z}$ also include every $z$ variable.
For every $q\ge0$,
\[
 F\in P_{A,z}^q\quad\Longleftrightarrow\quad
 F_\gamma\in P_A^{q-|\gamma|}\text{ for every }|\gamma|<q.
\]
\end{lemma}
\begin{proof}
Split each exponent uniquely as $(\beta,\gamma)$. Its weight for the
whole centre is the selected $w$ weight of $\beta$ plus $|\gamma|$.
The actual coordinate-power coefficient kernel proves the forward
implication coefficient by coefficient. Conversely, if that sum is less
than $q$, then $|\gamma|<q$ and the selected $w$ weight is less than
$q-|\gamma|$, so the assumed coefficient-series power makes the literal
coefficient of $F$ zero. The same actual kernel gives the reverse whole
ideal inclusion. No convergence, chosen quotient or finite truncation is
assumed.
\end{proof}

\begin{lemma}[Actual functional powers give the original whole plane-germ powers]
\label{hir:actual-original-plane-functional-power-germ-membership}
\uses{hir:relative-normal-entire-plane-kernel}
\lean{Hironaka.HironakaActualCoefficientPresentationCentrePower.mk_mem_normalIdeal_pow_of_functional_power,Hironaka.HironakaActualCoefficientPresentationCentrePower.actualUnshiftedTaylor_coordinate,Hironaka.HironakaActualCoefficientPresentationCentrePower.originalTaylor_normalCoordinate,Hironaka.HironakaActualCoefficientPresentationCentrePower.map_normalIdeal_originalTaylor,Hironaka.HironakaActualCoefficientPresentationCentrePower.originalTaylor_mem_coordinate_power_of_functional_power}
\leanok
Let $C=\{y=0\}\subset\mathbb C^{k+s}$. If a holomorphic representative
$f$ belongs at the origin to the genuine functional ideal
$\mathcal I_C^q$, then its original holomorphic germ belongs to the
$q$-th power of the actual ideal generated by the same $y$ coordinate
germs. Taylor expansion takes this ideal power into the actual formal
ideal power of those same variables.
\end{lemma}
\begin{proof}
For $q=0$ the ring ideal is the unit ideal. For a successor, the
functional power is the actual finite holomorphic span of products of a
previous-power member and a member vanishing on $C$. Its latter germ
belongs to the true coordinate-germ ideal by the already proved entire
plane-kernel dictionary. Induction, the actual product identity for
holomorphic germs, and their actual finite-sum identity place the full
representative in the ring ideal power. Taylor expansion is the genuine
original algebra homomorphism; its value on every original coordinate
is that same formal variable, so ideal-map and ideal-power operations
supply the final inclusion.
\end{proof}

\begin{theorem}[A genuine coefficient-marked centre internally supplies the full formal bound]
\label{hir:actual-coefficient-presentation-cosupport-centre-power}
\uses{hir:actual-formal-centre-power-relative-coefficients,hir:actual-original-plane-functional-power-germ-membership,hir:hadamard-bridge}
\source{bierstone-milman97-presentation-window-private:page-0071,bierstone-milman97-presentation-window-private:page-0072}
\lean{Hironaka.HironakaActualCoefficientPresentationCentrePower.mem_actualPresentationCentre_power_of_original_coefficient_cosupport}
\leanok
Let $F=\sum_\gamma\widehat{c_\gamma}(u,y)z^\gamma$ have actual original
holomorphic coefficient representatives $c_\gamma$. For each
$|\gamma|<q$, let $\mathfrak a_\gamma$ be an original coherent ideal
containing $c_\gamma$ at the origin. Assume the actual smooth coordinate
centre $C=\{y=0\}$ lies in the original coefficient cosupport
\[
 C\subset\operatorname{cosupp}(\mathfrak a_\gamma,q-|\gamma|)
 \qquad (|\gamma|<q).
\]
Then $F$ belongs to the actual ambient centre ideal $(y,z)^q$.
In particular this constructs, rather than assumes, the required centre
power of every original vertex member represented by those coefficients.
\end{theorem}
\begin{proof}
For each low coefficient the original coherent Hadamard bridge applies
to the genuine smooth centre and its actual cosupport containment. It
puts the coefficient representative in the functional power
$\mathcal I_C^{q-|\gamma|}$ at the origin. The previous lemma puts its
original germ, and then its genuine Taylor series, in the corresponding
coordinate ideal power. The coefficient identity in the original data
and the first lemma now put the ENTIRE formal member in $(y,z)^q$.
This argument does not assume that arbitrary monic representatives have
the needed centre order, nor that the formal properties (7.2)(4),(5) or
the full semicoherent presentation have already been constructed.
\end{proof}
